\documentclass[a4paper,11pt]{article}
\usepackage[margin=27mm]{geometry}
\usepackage{fontspec,amsmath,amssymb,booktabs,longtable,hyperref}
\setmainfont{STIXGeneral}
\setlength{\emergencystretch}{2em}
\hypersetup{colorlinks=true,urlcolor=blue,linkcolor=blue}
\title{Topic D: historical and current proof dependencies}
\author{Human Compression project}
\date{11 October 2026}
\begin{document}
\maketitle
This packet records the dependencies of the accepted Topic D edition. Its
three source contracts concern commuting involutions, Suzuki regular
polyhedra, and the Hall--Janko geometry census. The graph preserves all
37 source records and all 19 bound statement/proof contexts. It does not
assign proof weights, calculate compression factors, or add a source result.

An arrow labelled \emph{proof use}, \emph{definition use}, \emph{hypothesis
use}, \emph{proof assembly}, or \emph{retained input use} points from a
prerequisite to the argument that uses it. The other arrow types record
attribution, comparison, diagnostics, formal evidence or execution. They are
excluded from logical closures. Historical arrows document the dependencies
asserted or used in the source; they do not certify a flawed intermediate
argument retrospectively.

\section*{The three current routes}
\begin{description}
\item[Commuting involutions (L066).]
The orbit data of a marked Klein four-group gives five multiplicities,
faithfulness and parity. Reversible coordinates yield the ordered counts.
Endpoint exchange is then counted with its fixed-orbit correction. Ordinary
symmetric/alternating automorphism theorems remain retained inputs.
Degree six is handled by full automorphism orbits, middle-involution
stabilizers, generation and every intersection test. The complete exceptional
program is a necessary premise of this finite classification. The
regular-polytope correspondence is retained when interpreting the two
symmetric-group survivors and the absence of an alternating-group survivor.

\item[Suzuki polyhedra (L074).]
Classical finite Suzuki root groups, ovoids, automorphisms, canonical
subfields and maximal subgroups remain visible premises. The application
proves the coordinate normalization, positions a containing subfield,
identifies the generated group, checks the intersection property and proves
the rank restriction. These local steps identify actual field support;
the common support theorem does not supply them for free. Möbius inversion
and Frobenius orbit lengths then give the uniform count. The odd-orbit
lemma justifies dividing by two for duality. The stabilizer-fiber theorem
also gives the nested formula printed in the source. Finite field and
$q=8$ diagnostics are not premises of these all-parameter proofs.

\item[Hall--Janko geometries (L104).]
The complete subgroup catalogue, exact group operations and full
automorphism group remain retained. Coset realization, the typed Helly
criterion and arbitrary-Borel panel recognition connect subgroup tuples
to the exact geometric axioms. Transporters and stabilizer fibers make
maximal-residue augmentation complete. Four required programs perform
augmentation, full automorphism/type fusion and separate final geometry
verification using the common predicates. The resulting census is
$8,19,3$ in ranks two, three and four; the empty terminal level excludes
higher ranks. Catalogue completeness is not inferred from a successful
search.
\end{description}

\section*{What is shared, and what remains local}
The common extension theorem classifies data above one retained base orbit
by the action of its stabilizer. It is used for exceptional commuting
triples, Suzuki first-coordinate fibers, and the augmentation search.
The combined support theorem adds invariant admissibility, an equivariant
support poset and finite fixed-point counts. Its applications still owe
their admissibility, generation, support and freeness checks.

The finite-action centralizer calculation and the panel-recognition
arguments serve different branches. The full faithful-action normalizer
exact sequence is proved in the common paper, but the selected application
routes do not require that full statement. It remains present in the
complete paper; the graph does not silently delete its written cost.
No splitting of that sequence is asserted.

\section*{Historical qualifications}
The L066 endpoint-reversal statement is recovered with a repaired proof:
conjugating one endpoint alone does not establish simultaneous
conjugation of the pair. The source's self-duality remark is retained only
at the endpoint level. Its final alternating-group exclusion remains an
exclusion of \emph{regular} polyhedra. Earlier centralizer and polyhedron
results keep their original attribution.

L074 explicitly quotes earlier Suzuki and Lüneburg results. The current
paper retains their specified finite structure. It does not rederive the
full arbitrary-field Bruhat statement or the general ovoid plane-section
assertion. The printed torus-order coprimality lemma is unused by the new
intersection proof. Fermat's congruence is classical background. The old
almost-simple rank theorem is an input to the historical proof but is
rederived in the current application; this earns no separate source credit.

The L104 principal census does not include complete recovery of the printed
rank-four diagrams, all automorphism/correlation groups, the doubling
relation or every supplement annotation. Those assertions have historical
nodes and explicit excluded routes. The old Magma/geometry programs are
not identified with the present GAP implementation or TomLib catalogue.
Their complete construction and software costs remain unknown.

\section*{Programs, Lean and the review boundary}
The graph retains five complete required mathematical programs, three
complete diagnostics, and two complete execution drivers. Diagnostics have
validation arrows; drivers have operational arrows. The driver's historical
M22 branch is preserved in the actual file but is not a Topic D theorem
route. No program is shortened for this graph.

The original two Lean modules supply ten scoped reports: eight concern
fiber actions and odd-period arguments, and two give an abstract rank
obstruction. The later support supplement supplies 28 reports, comprising
the unchanged first twelve and sixteen further reports. The first twelve
are not added again. These are evidence for precisely scoped common
statements, not 38 independent mathematical results or formal proofs of
every application. Source-specific generation, the finite-field support
adapter, the number-theoretic divisor bridge and catalogue completeness
are not certified by that supplement.

Exact locators bind existing TeX, source PDFs and locally available page
extractions. They may overlap and do not constitute exclusive text ownership.
The portable checker verifies these bytes, the 37-row boundary, arrow types,
acyclicity, closures and required program routes. Corruption controls test
false formal premises, missing code, erased qualifications and invented
recovery of excluded assertions. Independent mathematical reading is still
required to accept the chosen edges. No source text, PDF, font or raster is
redistributed with this packet.
\end{document}
